\documentclass[12pt]{article}
\usepackage[utf8]{inputenc}
\usepackage[T1]{fontenc}
\usepackage[margin=1in]{geometry}
\usepackage{setspace}
\usepackage{xcolor}
\usepackage[colorlinks=true, urlcolor=blue!60!black, linkcolor=blue!60!black]{hyperref}
\setlength{\parindent}{2em}
\setlength{\parskip}{1.25em}
\renewcommand{\baselinestretch}{1.0}
\title{\textit{Myers v. United States}, 272 U.S. 52 (1926)}
\author{Case Brief}
\date{}

\begin{document}
\maketitle

\section{Case Metadata}
\begin{description}
    \item[Citation:] \textit{Myers v.\ United States}, 272 U.S.\ 52 (1926) (No.\ 2).
    \item[Decided:] October 25, 1926 (Argued December 5, 1924; Reargued April 13--14, 1925).
    \item[Vote:] 6--3.
    \item[Author:] Chief Justice William Howard Taft (joined by Justices Van Devanter, Sutherland, Butler, Sanford, and Stone).
    \item[Concurrences/Dissents:] No concurring opinions. Separate dissenting opinions were filed by Justice James Clark McReynolds (272 U.S.\ at 178), Justice Louis D.\ Brandeis (272 U.S.\ at 240), and Justice Oliver Wendell Holmes Jr.\ (272 U.S.\ at 177).
    \item[Chapter/Topic:] Separation of Powers and the Executive (Article II Vesting and Take Care Clauses; Presidential Removal Power; Unitary Executive Theory).
    \item[Primary Source URL:] \url{https://www.law.cornell.edu/supremecourt/text/272/52} (Justia: \url{https://supreme.justia.com/cases/federal/us/272/52/}; \textit{Trump v.\ Slaughter} (2026): \url{https://supreme.justia.com/cases/federal/us/609/25-332/}).
\end{description}

\section{Facts and Procedural History}
On July 21, 1917, President Woodrow Wilson appointed Frank S.\ Myers, by and with the advice and consent of the Senate, to a four-year term as Postmaster of the First Class in Portland, Oregon. Myers's appointment was governed by Section 6 of the Act of July 12, 1876, 19 Stat.\ 80, which provided that ``[p]ostmasters of the first, second, and third classes shall be appointed and may be removed by the President by and with the advice and consent of the Senate, and shall hold their offices for four years unless sooner removed or suspended according to law.''

On January 20, 1920, the Post Office Department demanded Myers's resignation. Myers refused. On February 2, 1920, he was removed from office by order of the Postmaster General acting at President Wilson's direction, without obtaining or seeking the Senate's advice and consent. Myers promptly petitioned the President and the Senate Committee on Post Offices for a hearing and protested his removal continuously until his four-year term expired on July 21, 1921, taking no other employment in the interim.

On April 21, 1921, Myers sued the United States in the Court of Claims to recover his unpaid salary from the date of removal through the end of his term (\$8,838.71). After Myers died, his administratrix (Lois P.\ Myers) continued the litigation. The Court of Claims entered judgment for the United States on the ground of laches (delay in filing suit). On direct appeal, the Supreme Court rejected the laches defense (which the Solicitor General effectively conceded at argument), restored the case for reargument, and invited Senator George Wharton Pepper of Pennsylvania to argue as \textit{amicus curiae} in defense of Congress's statutory restriction.

\section{Issue(s) Presented}
Does Article II of the Constitution vest in the President the exclusive power to remove executive officers appointed by him by and with the advice and consent of the Senate, rendering unconstitutional a federal statute that conditions such removal on the Senate's advice and consent?

\section{Holding}
\textbf{Yes.} Article II grants the President the exclusive constitutional power to remove executive officers appointed by him with the advice and consent of the Senate, and Congress may not condition or restrict that removal power by requiring Senate concurrence. Section 6 of the Act of July 12, 1876 (as well as the Reconstruction-era Tenure of Office Act of 1867) is unconstitutional. Judgment of the Court of Claims affirmed on constitutional grounds.

\section{Reasoning}
In a monumental 71-page opinion drawing directly on his own experience as the nation's 27th President, Chief Justice Taft marshaled constitutional text, structural logic, and early historical practice to articulate the foundational blueprint of the unitary executive.

\textbf{1. The Article II Vesting and Take Care Clauses.}
Taft contrasted Article I, \S~1 (``All legislative Powers \emph{herein granted} shall be vested in a Congress'') with Article II, \S~1 (``The executive Power shall be vested in a President of the United States of America''). Adopting Alexander Hamilton's \textit{Pacificus No.\ 1} construction, Taft reasoned that Article II's Vesting Clause is a general grant of executive power, subject only to express constitutional exceptions. Because the President cannot execute the laws ``alone and unaided,'' he must rely on subordinate executive officers; and because Article II, \S~3 commands him to ``take Care that the Laws be faithfully executed,'' the President must have ``general administrative control of those executing the laws,'' including the unrestricted power to remove subordinates ``for whom he cannot continue to be responsible.''

\textbf{2. Removal as Incident to Appointment---Without Extending the Senate's Check.}
Taft affirmed the settled principle that the power of removal is incident to the power of appointment. Crucially, however, he held that the Senate's role under Article II, \S~2 in providing ``Advice and Consent'' to appointments is a narrow, express exception to the separation of powers---designed at the Convention to reassure smaller states---that must be strictly construed and never extended by implication to removals. Blocking a nominee merely requires the President to pick another qualified candidate; blocking a removal forces the President to retain a disloyal, incompetent, or obstructive subordinate, paralyzing executive accountability.

\textbf{3. The ``Decision of 1789'' and 74 Years of Constitutional Liquidation.}
Taft devoted extensive historical analysis to the First Congress's 1789 debate over establishing the Department of Foreign Affairs. Led by James Madison, the House deliberately amended the bill (via Egbert Benson's amendment) to imply that the President's removal power flowed directly from the Constitution rather than from a revocable legislative grant, and the Senate concurred via Vice President John Adams's tie-breaking vote. Taft emphasized that this ``contemporaneous legislative exposition'' by Framers serving in the First Congress was acquiesced in by all three branches for 74 years (1789--1863) and recognized in \textit{Ex parte Hennen}, 13 Pet.\ 230 (1839), and \textit{Parsons v.\ United States}, 167 U.S.\ 324 (1897). He dismissed Chief Justice Marshall's contrary statement in \textit{Marbury v.\ Madison}, 1 Cranch 137 (1803), as \textit{obiter dictum} (noting Marshall himself later endorsed the Decision of 1789 in his \textit{Life of Washington}). Taft treated the Civil War/Reconstruction-era Tenure of Office Act of 1867 (and the 1876 postmaster statute) as partisan aberrations repeatedly resisted by Presidents Johnson, Grant, Cleveland, Wilson, and Coolidge.

\textbf{4. Inferior Officers and the Excepting Clause.}
Myers's counsel argued that because a postmaster is an ``inferior Officer'' whom Congress could have vested for appointment in the Postmaster General under Article II, \S~2's Excepting Clause (and whose removal Congress could then regulate under \textit{United States v.\ Perkins}, 116 U.S.\ 483 (1886)), Congress could likewise restrict the President's removal of a postmaster. Taft rejected this ``greater-includes-the-lesser'' argument: Congress's authority to regulate the removal of inferior officers is strictly conditioned on vesting their appointment in someone \emph{other} than the President with Senate consent, and in no event may Congress ``draw to itself, or to either branch of it, the power to remove or the right to participate in the exercise of that power.''

\textbf{5. Sweeping Reach to Quasi-Judicial Duties.}
In sweeping language that would ignite a century-long battle over independent agencies, Taft declared that even when executive officers or tribunals perform ``duties of a quasi judicial character,'' although the President cannot dictate a particular decision, he ``may consider the decision after its rendition as a reason for removing the officer'' to fulfill his Take Care duty.

\section{Concurring and Dissenting Opinions}
All three dissenters wrote separately, forming an unusual coalition spanning the Court's conservative and progressive wings.

\textbf{Justice Holmes, dissenting} (272 U.S.\ at 177): In a terse, two-paragraph dissent, Holmes dismissed Taft's structural deductions as ``spider's webs inadequate to control the dominant facts.'' Because Congress creates the office of postmaster, fixes its duration and pay, and may abolish it tomorrow, Congress can attach reasonable conditions to removal. The President's duty to take care that the laws be faithfully executed ``means no more than that the President must enforce the acts of Congress''---including the 1876 Act's removal restriction.

\textbf{Justice McReynolds, dissenting} (272 U.S.\ at 178): In an exhaustive dissent, McReynolds argued that the Constitution contains no express grant of an illimitable removal power and that inferring one from the general phrase ``executive Power'' violates the principle of enumerated powers. Distinguishing execution from legislation, he maintained that while the actual ouster of an officer is executive, prescribing the statutory conditions of tenure for an inferior office created by Congress is legislative. He also disputed Taft's reading of the 1789 debates, noting that only a minority of House speakers endorsed an indefeasible constitutional removal power and warning that Taft's broad theory threatened the independence of the ICC, FTC, Federal Reserve Board, and statutory courts.

\textbf{Justice Brandeis, dissenting} (272 U.S.\ at 240): Reviewing 137 years of federal statutes and state constitutional practice, Brandeis demonstrated that Congress had repeatedly restricted the removal of inferior officers with presidential approval and that the Decision of 1789 addressed only the removal of high-level department heads serving at pleasure, not inferior officers appointed for fixed terms. Rejecting Taft's claim that executive efficiency demanded unchecked presidential firing power, Brandeis delivered his celebrated defense of constitutional checks and balances: the separation of powers was adopted ``not to promote efficiency but to preclude the exercise of arbitrary power.''

\section{Key Quotations}
\begin{quote}
``The vesting of the executive power in the President was essentially a grant of the power to execute the laws. But the President alone and unaided could not execute the laws. He must execute them by the assistance of subordinates. ... As he is charged specifically to take care that they be faithfully executed, the reasonable implication, even in the absence of express words, was that as part of his executive power he should select those who were to act for him under his direction in the execution of the laws. The further implication must be, in the absence of any express limitation respecting removals, that as his selection of administrative officers is essential to the execution of the laws by him, so must be his power of removing those for whom he cannot continue to be responsible.'' (272 U.S.\ at 117)
\end{quote}

\begin{quote}
``Our conclusion on the merits, sustained by the arguments before stated, is that article 2 grants to the President the executive power of the government---i.e., the general administrative control of those executing the laws, including the power of appointment and removal of executive officers---a conclusion confirmed by his obligation to take care that the laws be faithfully executed; ... and, finally, that to hold otherwise would make it impossible for the President, in case of political or other difference with the Senate or Congress, to take care that the laws be faithfully executed.'' (272 U.S.\ at 163--164)
\end{quote}

\begin{quote}
``Then there may be duties of a quasi judicial character imposed on executive officers and members of executive tribunals whose decisions after hearing affect interests of individuals, the discharge of which the President cannot in a particular case properly influence or control. But even in such a case he may consider the decision after its rendition as a reason for removing the officer, on the ground that the discretion regularly entrusted to that officer by statute has not been on the whole intelligently or wisely exercised. Otherwise he does not discharge his own constitutional duty of seeing that the laws be faithfully executed.'' (272 U.S.\ at 135)
\end{quote}

\begin{quote}
``Generally, the actual ouster of an officer is executive action; but to prescribe the conditions under which this may be done is legislative. The act of hanging a criminal is executive; but to say when and where and how he shall be hanged is clearly legislative.'' (272 U.S.\ at 184, McReynolds, J., dissenting)
\end{quote}

\begin{quote}
``The doctrine of the separation of powers was adopted by the convention of 1787, not to promote efficiency but to preclude the exercise of arbitrary power. The purpose was not to avoid friction, but, by means of the inevitable friction incident to the distribution of the governmental powers among three departments, to save the people from autocracy.'' (272 U.S.\ at 293, Brandeis, J., dissenting)
\end{quote}

\section{Doctrinal Significance and Modern Legacy (through 2026)}
Few Supreme Court opinions have experienced as dramatic a century-long doctrinal arc---from sweeping triumph, to near-exile, to full centennial vindication---as \textit{Myers v.\ United States}.

\textbf{1. Retrenchment in the New Deal and Mid-Twentieth Century (1935--1988).}
Just nine years after \textit{Myers}, a unanimous Court sharply curtailed Taft's opinion in \textit{Humphrey's Executor v.\ United States}, 295 U.S.\ 602 (1935), upholding statutory ``for-cause'' removal protections (``inefficiency, neglect of duty, or malfeasance in office'') for Federal Trade Commission (FTC) Commissioners. Writing for the Court, Justice Sutherland (who had joined Taft's majority in \textit{Myers}) confined \textit{Myers} to ``purely executive officers'' like postmasters and disapproved its broader statements as \textit{dicta}, characterizing the FTC as a non-executive, ``quasi-legislative'' and ``quasi-judicial'' body. Subsequent decisions built on this functionalist turn: \textit{Wiener v.\ United States}, 357 U.S.\ 349 (1958) (implying for-cause protection for the adjudicatory War Claims Commission), and \textit{Morrison v.\ Olson}, 487 U.S.\ 654 (1988) (upholding for-cause protection for an inferior-officer independent counsel under a functional ``undue interference'' test, over Justice Scalia's famous lone dissent). Even during this era, however, \textit{Myers}'s core prohibition against \emph{congressional participation} in removal remained bedrock law, reaffirmed in \textit{Bowsher v.\ Synar}, 478 U.S.\ 714 (1986).

\textbf{2. The Roberts Court's Revival of \textit{Myers} (2010--2020).}
Beginning in \textit{Free Enterprise Fund v.\ Public Company Accounting Oversight Board}, 561 U.S.\ 477 (2010) (invalidating dual-layer for-cause removal protections for PCAOB members), and accelerating in \textit{Seila Law LLC v.\ Consumer Financial Protection Bureau}, 591 U.S.\ 197 (2020) (invalidating for-cause removal protection for the single director of the CFPB), the Roberts Court restored \textit{Myers} as the baseline constitutional rule: the President's ``general administrative control'' includes an unrestricted default removal power, and \textit{Humphrey's Executor} was treated as a narrow, historically anomalous exception not to be extended to new agency designs.

\textbf{3. Centennial Vindication: \textit{Trump v.\ Slaughter}, 609 U.S.\ \_\_\_ (2026).}
Exactly one century after \textit{Myers}, the Supreme Court completed its full doctrinal vindication in \textit{Trump v.\ Slaughter}, 609 U.S.\ \_\_\_ (No.\ 25--332, decided June 29, 2026). After President Donald J.\ Trump dismissed Democratic FTC Commissioners Rebecca Slaughter and Alvaro Bedoya in March 2025 without alleging statutory cause, the Court reversed the D.C.\ District Court in a 6--3 decision authored by Chief Justice Roberts (joined in full by Justices Alito, Gorsuch, Kavanaugh, and Barrett, and joined by Justice Thomas as to all but Part III--B). Opening with an explicit tribute---``One hundred years ago, this Court honored the Convention's choice in the seminal case of \textit{Myers v.\ United States}, 272 U.S.\ 52 (1926)''---the Court formally overruled \textit{Humphrey's Executor} and invalidated the FTC's for-cause removal statute (15 U.S.C.\ \S~41). Holding that the FTC's sweeping rulemaking, in-house adjudication, and federal-court enforcement powers are ``quintessentially executive,'' Roberts reaffirmed that all principal officers exercising executive power must be removable by the President at will to preserve democratic accountability through Madison's ``chain of dependence.''

Three nuances define the post-\textit{Slaughter} landscape in 2026:
\begin{itemize}
    \item \textbf{Narrow Historical Carve-Outs:} In Part III--B (which Justice Thomas declined to join) and in the companion decision \textit{Trump v.\ Cook}, 609 U.S.\ \_\_\_ (2026), the Court preserved a narrow exception for the Federal Reserve, grounded in the distinct historical tradition of the First and Second Banks of the United States, while expressly reserving the status of tenure protections for non-Article III courts (such as the Tax Court and Court of Federal Claims).
    \item \textbf{Justice Gorsuch's Concurrence and the Delegation Reckoning:} Concurring in \textit{Slaughter}, Justice Gorsuch celebrated \textit{Myers}'s vindication but warned that eliminating independent-agency insulation transfers vast delegated legislative and judicial powers directly into the President's hands (``The ability to judge disputes in-house remains, but now the house is white''), urging the Court to pair \textit{Myers} with rigorous enforcement of the nondelegation doctrine, the major questions doctrine, and Article III/Seventh Amendment limits (\textit{SEC v.\ Jarkesy}, 603 U.S.\ 109 (2024)).
    \item \textbf{Dissenting Critique:} Justice Sotomayor (joined by Justices Kagan and Jackson) dissented, attacking both \textit{Myers}'s historical account of the Decision of 1789 and the majority's abandonment of 91 years of \textit{stare decisis} and congressional reliance across dozens of multimember agencies (such as the NLRB, FERC, CPSC, NRC, and MSPB).
\end{itemize}

\section{Pedagogical Critiques and Exam Hypotheticals}
\textbf{Pedagogical Critiques:}
\begin{enumerate}
    \item \textbf{Originalist Contestability of the Decision of 1789:} Modern legal historians (cited extensively by Justice Sotomayor in her \textit{Slaughter} dissent) argue that Taft overstated the consensus of the First Congress: the 1789 House was fractured into three or four competing factions, and its vote rejecting \emph{Senate participation} in removal did not necessarily resolve whether Congress could impose \emph{for-cause} tenure limits via statute.
    \item \textbf{The Ratchet Effect of Unitary Removal Without Nondelegation:} As Justice Gorsuch highlighted in \textit{Trump v.\ Slaughter} (2026), combining \textit{Myers}'s plenary presidential removal power with twentieth-century broad delegations of rulemaking and adjudicative authority concentrates unprecedented power in a single executive---creating an institutional tension that neither Taft in 1926 nor the \textit{Slaughter} majority fully resolved.
\end{enumerate}

\textbf{``Opposite Side'' Contrast (\textit{Myers} and \textit{Trump v.\ Slaughter} vs.\ \textit{Trump v.\ Cook} (2026) and \textit{Morrison v.\ Olson} (1988)):}
In \textit{Myers} (1926) and \textit{Trump v.\ Slaughter} (2026), the Court held that whenever an agency exercises executive law-enforcement, rulemaking, or administrative power (whether a single postmaster or a five-member FTC), Article II forbids Congress from insulating its principal officers from at-will presidential removal, whether by requiring Senate consent (\textit{Myers}) or by imposing statutory ``for-cause'' restrictions (\textit{Slaughter}). By contrast, \textit{Morrison v.\ Olson}, 487 U.S.\ 654 (1988) (whose continued vitality after \textit{Slaughter} remains sharply contested for inferior officers) upheld a modest for-cause restriction on an \emph{inferior} officer with limited, temporary jurisdiction appointed by a court of law; and \textit{Trump v.\ Cook}, 609 U.S.\ \_\_\_ (2026), preserved removal insulation for Federal Reserve Governors based on a uniquely entrenched founding-era historical tradition of central-bank independence (the First and Second Banks of the United States) that does not extend to ordinary regulatory commissions.

\textbf{Exam Hypothetical:}
Congress enacts the Artificial Intelligence Safety Act of 2026, establishing the five-member Federal AI Commission (FAIC), appointed by the President with Senate confirmation for staggered seven-year terms, with no more than three members from the same political party. The FAIC is empowered to promulgate binding safety rules for frontier AI models, investigate technology firms, adjudicate violations in-house, and seek billion-dollar civil penalties in federal court. The statute provides two removal clauses:
\begin{itemize}
    \item \textbf{Clause 1:} ``Commissioners may be removed by the President only for inefficiency, neglect of duty, or malfeasance in office.''
    \item \textbf{Clause 2:} ``The FAIC's Chief Inspector---an inferior officer appointed by the FAIC Chair---may be removed by the Chair only upon a showing of good cause.''
\end{itemize}
The President fires two FAIC Commissioners solely for policy disagreements and orders the new Chair to fire the Chief Inspector at will. Both the Commissioners and the Chief Inspector sue.
\textit{Analysis:} Under \textit{Myers v.\ United States} as fully vindicated by \textit{Trump v.\ Slaughter} (2026), \textbf{Clause 1 is plainly unconstitutional}: the FAIC wields quintessential executive powers (binding rulemaking, investigations, in-house adjudication, and civil-penalty litigation), so its principal officers must be removable by the President at will, and \textit{Humphrey's Executor} is no longer good law. Nor can the FAIC claim the \textit{Trump v.\ Cook} historical carve-out applicable to the Federal Reserve. As to \textbf{Clause 2}, the Chief Inspector is an inferior officer whose appointment Congress vested in a Head of Department (the Commission/Chair) under Article II, \S~2's Excepting Clause. Under \textit{Myers}'s own discussion of \textit{United States v.\ Perkins}, 116 U.S.\ 483 (1886) (and \textit{Morrison v.\ Olson}), Congress may impose a \emph{single} layer of for-cause protection on an inferior officer appointed by a department head, provided the department head is removable at will by the President (avoiding the unconstitutional \emph{dual-layer} tenure problem of \textit{Free Enterprise Fund}). Once Clause 1 is severed under \textit{Slaughter}, the Chair is removable at will, meaning Clause 2 creates only a single layer of protection---though challengers could still argue under a maximalist reading of \textit{Myers} and \textit{Slaughter} that even single-layer limits on inferior officers exercising core enforcement power impair the President's Take Care duty.

\end{document}
